\documentclass[12pt]{book}
\usepackage{amsmath}
\usepackage{amssymb}
\usepackage{geometry}
\geometry{margin=1in}

\title{Volume 05: Nuclear Physics and Particle Architecture \\ Book 01: Nucleon Structures}
\author{James C. Harvey}
\date{December 2025}

\begin{document}
\maketitle
\tableofcontents

\chapter{Proton Geometry}

\section*{Abstract}
The proton is a stable toroidal boundary with trefoil topology. Its magnetic moment and coupling behavior
follow from curvature, circulation, and coupling efficiency.

\section{Trefoil Structure}
The proton is a $6\pi$ trefoil torus. Its curvature is
\begin{equation}
\kappa_P = \frac{1}{R_P}.
\end{equation}
The circulation factor $\Gamma_P$ and slip $\eta_P$ determine coupling.

\section{Curvature Scaling}
If $R_P$ changes by $\Delta R_P$, then
\begin{equation}
\Delta \kappa_P = -\frac{\Delta R_P}{R_P^2},
\end{equation}
so curvature increases with compaction and directly enhances coupling.

\section{Proton Moment}
\begin{equation}
\mu_P \propto \Gamma_P \kappa_P (1-\eta_P).
\end{equation}

\chapter{Proton Stability}

\section*{Abstract}
Stability is the alignment of circulation with boundary curvature. This chapter formalizes the stability
condition in SDT terms.

\section{Stability Condition}
For a stable proton,
\begin{equation}
\dot{E}_P = P_{\text{CMB}} A_{\text{eff}} \Gamma_P \kappa_P (1-\eta_P)
\end{equation}
is constant over time. Deviations correspond to excitation or decay regimes.

\chapter{Neutron Genesis}

\section*{Abstract}
The neutron is a composite structure: a proton boundary with a nested electron. Its negative magnetic
moment is the signature of reversed circulation.

\section{Nested Electron}
The nested electron occupies the central region of the proton torus. Its circulation is reversed,
producing negative magnetic moment:
\begin{equation}
\mu_N \propto -\Gamma_{E} \kappa_{E} (1-\eta_N).
\end{equation}

\chapter{Deuteron and Light Nuclei}

\section*{Abstract}
Light nuclei are coaxial packings of toroidal boundaries. This chapter summarizes the SDT geometry of
deuteron and related structures.

\section{Deuteron Structure}
The deuteron is a coupled proton--neutron pair with shared flux:
\begin{equation}
\mu_D = f_{\text{damp}}(\mu_P + \mu_N),
\end{equation}
where $f_{\text{damp}}$ is the geometry-based damping factor.

\section{Damping Interpretation}
The damping factor is a coupling efficiency:
\begin{equation}
f_{\text{damp}} = 1 - \eta_D,
\end{equation}
so deuteron coupling reduces to a single slip-controlled reduction of the combined moments.

\section{Cross-Reference}
Full derivations and numeric values are provided in the benchmark documents and calculation scripts.
\end{document}
